\section*{Hoofdstuk \ref{ch:b1:non-inertial-frames} --- Niet-inertiaalstelsels; dynamica op aarde}

\begin{solution}{exo:b1:non-inertial-frames:1}
$\tan\alpha = a/g = 0.255$, $\alpha = 14^\circ$ (het lood hangt naar achteren);
$T = m\sqrt{g^2 + a^2} = 1.03\,mg$. Remmen: dezelfde hoek en spankracht, het lood
naar voren gekanteld.
\end{solution}

\begin{solution}{exo:b1:non-inertial-frames:2}
$\Omega = 2\pi/4.0 = \qty{1.57}{rad/s}$; $\Omega^2r = \qty{7.4}{m/s^2} = 0.75g$,
in het meedraaiende stelsel naar buiten gericht; het zitje (wrijving en
beugel) levert de werkelijke kracht naar binnen.
\end{solution}

\begin{solution}{exo:b1:non-inertial-frames:3}
$\Omega = 2\pi/86164 = \qty{7.29e-5}{rad/s}$; $\Omega^2R = \qty{0.034}{m/s^2}$,
$0.35\%$ van $g$. Wegzweven vergt $\Omega^2R = g$: $T = 2\pi\sqrt{R/g} =
\qty{84}{min}$ --- zeventien keer sneller dan nu.
\end{solution}

\begin{solution}{exo:b1:non-inertial-frames:4}
$a_C = 2\Omega v\sin\lambda = 2 \times 7.29\times10^{-5} \times 27.8 \times 0.707
= \qty{2.9e-3}{m/s^2}$; $F = \qty{1.1}{kN}$. Naar het noorden: naar het oosten geduwd,
op de rechter (oostelijke) rail; naar het oosten: naar het zuiden geduwd, opnieuw
de rechterrail.
\end{solution}

\begin{solution}{exo:b1:non-inertial-frames:5}
$z = \Omega^2r^2/2g$: \qty{2.0}{mm} bij \qty{2}{rad/s}, \qty{5.1}{cm} bij
\qty{10}{rad/s}. Een paraboloïde bundelt een evenwijdige bundel in één punt:
draaiend kwik is een volmaakte, goedkope, zichzelf polijstende parabolische spiegel
(die alleen niet gekanteld kan worden).
\end{solution}

\begin{solution}{exo:b1:non-inertial-frames:6}
De centrifugale versnelling $\Omega^2R\cos\lambda$ wijst van de
as af; haar component loodrecht op de lokale verticaal is $\Omega^2R\cos\lambda
\sin\lambda$, die $\vect g$ kantelt over $\delta \approx \Omega^2R\sin\lambda\cos\lambda/g
= 0.034 \times 0.5/9.81 = \qty{1.7e-3}{rad} = 6'$. Evenaar: $g$ verminderd met
$\Omega^2R = \qty{0.034}{m/s^2}$ door de rotatie (de afplatting doet er nog meer bij).
\end{solution}

\begin{solution}{exo:b1:non-inertial-frames:7}
$\ddot x_E = 2\Omega gt\cos\lambda$: $x_E = \tfrac13\Omega g\cos\lambda\,t^3$; $t =
\sqrt{2h/g} = \qty{4.5}{s}$: $x_E = \tfrac13 \times 7.29\times10^{-5} \times 9.81
\times 0.707 \times 92 = \qty{1.6}{cm}$. Naar het oosten: de bovenkant van de schacht
beweegt sneller naar het oosten (grotere straal) dan de onderkant; de steen houdt dat
overschot --- of, in het meedraaiende stelsel, de \omterm{thm:b1:non-inertial-frames:dynamics}{corioliskracht} op een neerwaartse
snelheid wijst naar het oosten.
\end{solution}

\begin{solution}{exo:b1:non-inertial-frames:8}
$T_F = 23.93/\sin\lambda$: Parijs \qty{31.8}{h}; pool \qty{23.9}{h}; evenaar
oneindig (geen precessie). Parijs: $360^\circ/31.8 = 11.3^\circ$ per uur.
\end{solution}

\begin{solution}{exo:b1:non-inertial-frames:9}
$a_C = 2\Omega v\sin\lambda = \qty{2.1e-3}{m/s^2}$, naar rechts. Lucht die
van alle kanten samenstroomt wordt naar rechts afgebogen en spiraalt dus
tegen de klok in naar binnen (van boven gezien): de noordelijke cycloon. Een afvoerputje
loopt in seconden leeg en een tornado ontstaat in minuten: de coriolisafbuiging over
zulke tijden is verwaarloosbaar tegen elke beginwerveling.
\end{solution}

\begin{solution}{exo:b1:non-inertial-frames:10}
In het stelsel van de auto kantelt de schijnbare zwaartekracht $\vect g - \vect a$ naar achteren.
De lucht, zwaarder, ``valt'' naar achteren; de ballon, lichter dan de lucht die hij
verdringt, wordt langs de schijnbare verticaal omhoog gedreven --- hij helt naar \emph{voren},
tegengesteld aan een schietlood.
\end{solution}

\begin{solution}{exo:b1:non-inertial-frames:11}
$\Omega^2R = g$: $\Omega = \qty{0.31}{rad/s}$, $T = \qty{20}{s}$. $a_C = 2\Omega v
= \qty{0.94}{m/s^2}$, ongeveer $0.1g$: loopt hij met de draairichting mee, dan voelt hij
zich $10\%$ zwaarder (de \omterm{thm:b1:non-inertial-frames:dynamics}{corioliskracht} wijst naar buiten), er tegenin
lichter. Een losgelaten bal: in het inertiaalstelsel gaat hij rechtdoor terwijl
de vloer omhoog kromt; over de valtijd $\sqrt{2h/g} = \qty{0.64}{s}$ is de
afwijking van de orde $\tfrac13\Omega\,g\,t^3 \cdot 2 \approx 2 \times 0.31 \times
9.81 \times 0.26/3 \approx \qty{0.5}{m}$ --- enkele decimeters, tegen de
draairichting in.
\end{solution}

\begin{solution}{exo:b1:non-inertial-frames:12}
$GM/(d - r)^2 - GM/d^2 \approx 2GMr/d^3$ (eerste orde in $r/d$). Maan:
$2 \times 6.67\times10^{-11} \times 7.35\times10^{22} \times 6.37\times10^6/
(3.84\times10^8)^3 = \qty{1.1e-6}{m/s^2}$; zon: $2 \times 1.33\times10^{20}
\times 6.37\times10^6/(1.50\times10^{11})^3 = \qty{5.0e-7}{m/s^2}$; verhouding
$0.46$: de getijden van de zon zijn bijna de helft van die van de maan, hoewel haar
aantrekking $180$ keer groter is --- getijden gaan als $M/d^3$.
\end{solution}

\begin{solution}{pb:b1:non-inertial-frames:1}
\textbf{1.} $T = 2\pi\sqrt{67/9.81} = \qty{16.4}{s}$; $x_0/\ell = 0.045$
($2.6^\circ$): klein.

\textbf{2.} $v_{\max} = x_0\omega_0 = 3.0 \times 0.383 = \qty{1.15}{m/s}$; hoogte
$x_0^2/2\ell = \qty{6.7}{cm}$.

\textbf{3.} $T = mg + mv_{\max}^2/\ell = mg(1 + 0.002)$: $1.002\,mg$.

\textbf{4.} $\tfrac12\rho C_xSv_{\max}^2 = 0.5 \times 1.2 \times 0.05 \times 1.3 =
\qty{0.04}{N}$, tegen $m\omega_0^2x_0 = \qty{12}{N}$: driehonderd keer
kleiner, en toch genoeg om de slinger in een paar uur tot stilstand te brengen.

\textbf{5.} Een lange draad geeft een lange periode, dus een traag gewicht en een
kleine relatieve weerstand, en maakt de precessie per zwaai zichtbaar; een
zwaar gewicht slaat veel energie op die de weerstand mag opeten --- uren zwaaien.

\textbf{6.} Het gewicht (verticaal) en de spankracht (in het vlak van de draad en
de verticaal, dus in het zwaaivlak): geen enkele kracht heeft een component
loodrecht op dat vlak.

\textbf{7.} Het vlak ligt vast tussen de sterren; het ijs draait er
oostwaarts met $\Omega$ onderdoor, dus ziet de waarnemer het vlak met de klok mee draaien,
terug bij het begin na één sterrendag, \qty{23.93}{h}.

\textbf{8.} Elke halve zwaai tekent een iets gedraaide koorde: een rozet;
$86164/8.2 \approx 10500$ halve zwaaien, elk $0.034^\circ$ verder rond.

\textbf{9.} $\vect\Omega = \Omega(0, \cos\lambda, \sin\lambda)$.

\textbf{10.} Met $\vect v' = (\dot x, \dot y, 0)$:
$\vect\Omega\wedge\vect v' = \Omega(-\sin\lambda\,\dot y,\ \sin\lambda\,\dot x,\ -\cos\lambda\,\dot x)$,
dus heeft $-2m\vect\Omega\wedge\vect v'$ de horizontale componenten
$2m\Omega\sin\lambda\,(\dot y, -\dot x)$: alleen $\Omega\sin\lambda$ komt erin voor.

\textbf{11.} $\ddot x = -\omega_0^2x + 2\omega_F\dot y$, $\ddot y = -\omega_0^2y -
2\omega_F\dot x$; $\omega_F = 7.29\times10^{-5} \times 0.753 = \qty{5.5e-5}{rad/s}
\ll \omega_0 = \qty{0.38}{rad/s}$.

\textbf{12.} $\ddot u = -\omega_0^2u + 2\omega_F(\dot y - j\dot x) = -\omega_0^2u -
2j\omega_F\dot u$.

\textbf{13.} $-\alpha^2 - 2\omega_F\alpha + \omega_0^2 = 0$: $\alpha = -\omega_F \pm
\sqrt{\omega_F^2 + \omega_0^2} \approx -\omega_F \pm \omega_0$.

\textbf{14.} De haakjes geven een vlakke trilling (een vaste lijn door de
oorsprong voor reële $A = B$); de factor $\eu^{-j\omega_Ft}$ draait het hele
patroon met de klok mee met $\omega_F$: het zwaaivlak preceseert met de klok mee met
$\Omega\sin\lambda$.

\textbf{15.} $T_F = 2\pi/\omega_F = \qty{1.14e5}{s} = \qty{31.8}{h}$; $11.3^\circ$
per uur.

\textbf{16.} $\sin\lambda = 0$: de lokale verticale component van $\vect\Omega$
verdwijnt; de horizontale \omterm{thm:b1:non-inertial-frames:dynamics}{corioliskracht} is nul voor een horizontale beweging.

\textbf{17.} $2m\omega_Fv_{\max} = 2 \times 28 \times 5.5\times10^{-5} \times 1.15 =
\qty{3.5}{mN}$; gewicht \qty{275}{N} ($10^{-5}$ daarvan), terugdrijvende kracht
$m\omega_0^2x_0 = \qty{12}{N}$: drieduizend keer kleiner.

\textbf{18.} $\omega_FT/2 = 5.5\times10^{-5} \times 8.2 = \qty{4.5e-4}{rad}$:
$x_0 \times 4.5\times10^{-4} = \qty{1.4}{mm}$ per halve zwaai.

\textbf{19.} Pinnetjes die een centimeter uit elkaar op de ring staan vallen om de $\sim 7$
halve zwaaien, ongeveer elke minuut.

\textbf{20.} Een duw met de hand geeft een zijwaartse snelheid: het gewicht
beschrijft dan een ellips, waarvan de eigen trage precessie (een gevolg van de
eindige \omterm{def:b1:wave-propagation:signal}{amplitude}) die van Foucault maskeert.

\textbf{21.} Een radiale duw voegt energie toe langs de zwaai zonder
transversale snelheid toe te voegen, zodat het vlak onaangeroerd blijft; een zijwaartse
component zou de slinger in een ellips drijven en de precessie namaken of
verbergen.

\textbf{22.} Het precessietempo meet $\Omega\sin\lambda$, de verticale
component van de aardrotatie; kent men $\Omega$ uit de sterrenkunde, dan geeft het
tempo de breedte (en omgekeerd).

\textbf{23.} $\sin\lambda = 23.93/48 = 0.499$: $\lambda = 30^\circ$.

\textbf{24.} Zijn \omterm{def:b1:angular-momentum:L}{impulsmoment} blijft behouden, dus blijft zijn as vast
tussen de sterren terwijl de aarde eronder doordraait: het gyrokompas, dat
het ware noorden vindt zonder magneet.

\textbf{25.} De \omterm{thm:b1:non-inertial-frames:dynamics}{corioliskracht} $-2m\vect\Omega\wedge\vect v'$; haar horizontale
deel groeit als $\sin\lambda$; één omwenteling in $23.93\,\text{h}/\sin\lambda$,
\qty{31.8}{h} in Parijs.
\end{solution}
